
\thispagestyle{empty}
\begin{center}
\underline{INTRODUCTION}
\end{center}

Let $S$ be a scheme and $f:X\longrightarrow S$ a morphism of schemes. If $f$ is
proper and smooth, and if the prime number $\ell$ is invertible on $S$, the
$\ell$-adic cohomology groups $H^i(X_{\bar s},\mathbb Z_\ell)$ of the geometric
fibres of $f$ form an $\ell$-adic local system on $S$. When $f$ is only assumed
proper, these groups are the fibres of an $\ell$-adic sheaf
$R^i f_*\mathbb Z_\ell$ on $S$. The \underline{theory of vanishing cycles}
relates the ramification of this sheaf on $S$ to the singularities of $f$.

We shall consider only the case where $S$ is a henselian trait (= the spectrum
of a henselian discrete valuation ring). In practice, this is the essential
case. I refer to (I 2.1) for a heuristic description of the theory. Let
$S=\Spec(V)$, let $k(\bar\eta)$ be an algebraic closure of the fraction field
$k(\eta)$ of $V$, and let $k(\bar s)$ be the corresponding algebraic closure of
the residue field $k(s)$ ($k(\bar s)$ is the residue field of the normalization
$\bar V$ of $V$ in $k(\bar\eta)$). In the special case where $X$ is proper and
flat over $S$, of relative dimension $n$, and where $f$ has only one point of
non-smoothness $x\in X_s$, one defines $\Gal(\bar\eta/\eta)$-modules $\phi^i$,
of a purely local nature in a neighbourhood of $x$, vanishing for
$i\notin[0,n]$ (I 4.2), and one constructs a long exact sequence of
$\Gal(\bar\eta/\eta)$-modules

\[
\cdots \longrightarrow H^i(X_{\bar s},\mathbb Z_\ell)
\xrightarrow{\mathrm{sp}} H^i(X_{\bar\eta},\mathbb Z_\ell)
\longrightarrow \phi^i
\longrightarrow H^{i+1}(X_{\bar s},\mathbb Z_\ell)
\longrightarrow \cdots
\]

(the $\Gal(\bar s/s)$-module $H^i(X_{\bar s},\mathbb Z_\ell)$ is regarded as a
$\Gal(\bar\eta/\eta)$-module on which inertia acts trivially; $\mathrm{sp}$ is
the specialization map).

Criteria are also given, valid for the moment only in characteristic $0$, for
the $\phi^i$ with small $i$ to vanish (for example, if $X$ is moreover locally
a complete intersection, $\phi^i=0$ for $i\ne n$ (I. 4.5)).

\clearpage
\thispagestyle{empty}
\begin{center}
VI
\end{center}

The \underline{monodromy theorem} asserts that the action of the inertia
subgroup $I$ of $\Gal(\bar\eta/\eta)$ is quasi-unipotent: for $T\in I$, there
exist integers $N>0$, $M>0$ such that the endomorphism $(T^M-I)^N$ of
$H^i(X_{\bar\eta},\mathbb Z_\ell)$ is zero. Two proofs of this theorem are
given. The first (I.1.2), which is arithmetic in nature, applies as soon as the
cohomology groups under consideration have reasonable finiteness properties.
The key point is that, when $k(s)$ is of finite type over its prime subfield,
the action of $I$ is quasi-unipotent for every $\ell$-adic representation of
$\Gal(\bar\eta/\eta)$. The second proof, which is more geometric, requires
purity and resolution of singularities: for the moment it is valid only in
characteristic $0$ or for an $H^1$. It provides valuable information about the
nilpotence exponent $N$ ($N\leq i+1$ for an $H^i$). As will be seen later, over
$\mathbb C$ it is well suited to comparison with transcendental theory.

Applied to the $H^1$ of abelian varieties, i.e. to their Tate module, these
results make it possible to study their reduction modulo $p$. Thus two proofs
are given of the stable reduction theorem for abelian varieties, according to
which, after ramification, the connected special fibre of the Néron model is
an extension of an abelian variety by a torus. The first (I. 6) follows the
method of the arithmetic proof of the monodromy theorem. The second (IX 3.6)
rests on a much finer analysis of the Néron model and of its polarization
properties.

Grothendieck's Exposés I through V were not written up. They were summarized in
one Exposé I. The results stated there are proved succinctly, but essentially
completely.

Exposé II applies the method of Lefschetz pencils and (I 5.3) to the study of
the fundamental group. If $S$ is a surface over an algebraically closed field
$k$ of characteristic exponent $p$, it is shown that the profinite group
$\pi_1^{(p)}(S,s)$, the completion away from $p$ of the fundamental group,
admits a finite pro-$(p)$ presentation.

As explained above, Exposés III through V do not exist.

\clearpage
\thispagestyle{empty}
\begin{center}
VII
\end{center}

Exposé VI contains, with a few additions, Schlessinger's theory of
deformations. Care is taken there to account for the infinitesimal
automorphisms of the objects being classified, and not to pass too abruptly to
the functor of isomorphism classes of objects. It also gives a new construction
of $\underline{\operatorname{Ext}}^{i}(L_{X/S},-)$ ($L_{X/S}$ being the
relative cotangent complex of $X/S$). In the rest of the seminar this Exposé is
used chiefly through its study of deformations of ordinary quadratic
singularities.

Exposés VII and VIII are devoted to the theory of biextensions. This theory
plays an essential role in Exposé IX, in expressing what becomes of a
polarization when one passes from an abelian variety to its Néron model. Exposé
IX contains the proof of the stable reduction theorem for abelian varieties,
as well as various applications.

The continuation of this seminar, SGA 7 II, by P. Deligne and N. Katz, will
appear later.

\begin{flushright}
Bures-sur-Yvette, May 1972,\\
P. DELIGNE
\end{flushright}
